\documentclass[aspectratio=169,11pt]{beamer}
% Shared CMPSC 480 Beamer preamble.
\usepackage[T1]{fontenc}
\usepackage[utf8]{inputenc}
\usepackage{lmodern}
\usepackage{amsmath,amssymb,mathtools}
\usepackage{booktabs,array,tabularx}
\usepackage{xcolor}
\usepackage{listings}
\usepackage{tikz}
\usetikzlibrary{arrows.meta,positioning,shapes.geometric}

% Ignore only negligible TeX box diagnostics that are below visual tolerance.
% Larger layout defects still appear in the build log and in rendered QA.
\vfuzz=3pt
\hbadness=2000

\definecolor{courseink}{HTML}{111111}
\definecolor{coursegray}{HTML}{EDEDED}
\definecolor{courserule}{HTML}{B8BCC4}
\definecolor{courseblue}{HTML}{3D8DFF}
\definecolor{courselightblue}{HTML}{D0EDFA}
\definecolor{coursemuted}{HTML}{5C6470}
\definecolor{coursegreen}{HTML}{0B7A53}
\definecolor{coursered}{HTML}{B3261E}

\usetheme{default}
\usefonttheme{professionalfonts}
\setbeamertemplate{navigation symbols}{}
\setbeamersize{text margin left=0.65cm,text margin right=0.65cm}

\setbeamercolor{normal text}{fg=courseink,bg=white}
\setbeamercolor{frametitle}{fg=courseink,bg=white}
\setbeamercolor{title}{fg=courseink,bg=white}
\setbeamercolor{subtitle}{fg=coursemuted,bg=white}
\setbeamercolor{structure}{fg=courseblue}
\setbeamercolor{alerted text}{fg=coursered}
\setbeamercolor{block title}{fg=courseink,bg=courselightblue}
\setbeamercolor{block body}{fg=courseink,bg=coursegray}
\setbeamercolor{example text}{fg=coursegreen}

\setbeamerfont{title}{size=\fontsize{25}{29}\selectfont,series=\bfseries}
\setbeamerfont{subtitle}{size=\fontsize{13}{16}\selectfont}
\setbeamerfont{frametitle}{size=\fontsize{19}{22}\selectfont,series=\bfseries}
\setbeamerfont{normal text}{size=\fontsize{11}{14}\selectfont}
\setbeamerfont{footline}{size=\fontsize{7.5}{9}\selectfont}

\setbeamertemplate{frametitle}{%
  \nointerlineskip
  % Use content-driven height so a long assessment title can wrap without
  % being clipped by a fixed-height title bar.
  \begin{beamercolorbox}[wd=\paperwidth,sep=0.17cm,leftskip=0.48cm,rightskip=0.48cm]{frametitle}
    \insertframetitle
  \end{beamercolorbox}
  \vspace{-0.08cm}
  {\color{courserule}\rule{\textwidth}{0.35pt}}
}

\setbeamertemplate{footline}{%
  \leavevmode
  \begin{beamercolorbox}[wd=\paperwidth,ht=2.6ex,dp=1.1ex,leftskip=.65cm,rightskip=.65cm]{author in head/foot}
    \color{coursemuted}\insertshorttitle\hfill\insertframenumber/\inserttotalframenumber
  \end{beamercolorbox}
}

\setbeamertemplate{itemize item}{\color{courseblue}\small$\blacktriangleright$}
\setbeamertemplate{itemize subitem}{\color{courseblue}\scriptsize$\bullet$}
\setbeamertemplate{enumerate item}{\color{courseblue}\insertenumlabel.}

\lstdefinestyle{coursepython}{
  language=Python,
  basicstyle=\ttfamily\fontsize{8.5}{10}\selectfont,
  keywordstyle=\color{courseblue}\bfseries,
  commentstyle=\color{coursemuted}\itshape,
  stringstyle=\color{coursegreen},
  showstringspaces=false,
  columns=fullflexible,
  keepspaces=true,
  frame=single,
  rulecolor=\color{courserule},
  backgroundcolor=\color{coursegray},
  breaklines=true,
  tabsize=4,
  xleftmargin=0.15cm,
  xrightmargin=0.15cm,
  framexleftmargin=0.15cm,
  framexrightmargin=0.15cm
}
\lstset{style=coursepython}

\newcommand{\points}[1]{\hfill{\color{courseblue}\textbf{[#1 points]}}}
\newcommand{\deliverable}[1]{\textbf{\color{courseblue}#1}}
\newcommand{\code}[1]{\texttt{#1}}
\newcommand{\keyidea}[1]{%
  \begin{beamercolorbox}[rounded=false,sep=0.55em,wd=\linewidth]{block body}
    \textbf{Key idea.} #1
  \end{beamercolorbox}
}
\newcommand{\answerlabel}{\textbf{\color{coursegreen}Answer.}\ }
\newcommand{\gradinglabel}{\textbf{\color{courseblue}Grading guidance.}\ }

\newenvironment{questionblock}[1]
  {\begin{block}{#1}}
  {\end{block}}

\newenvironment{solutionblock}[1]
  {\begin{block}{#1}}
  {\end{block}}

\AtBeginSection[]{
  \begin{frame}
    \vfill
    {\usebeamerfont{title}\color{courseink}\insertsectionhead\par}
    \vspace{0.4cm}
    {\color{courseblue}\rule{0.32\paperwidth}{3pt}}
    \vfill
  \end{frame}
}


% CMPSC480_TOTAL_POINTS: 10
% CMPSC480_ITEM_IDS: 1,2,3,4
\title[Week 4 Assignment Solutions]{Week 4 Assignment: Planning in a Stochastic Gridworld}
\subtitle{MDP modeling, value iteration, policy iteration, and verification}
\author{CMPSC 480 - Instructor Solution Deck}
\date{}

\begin{document}


\begin{frame}
  \titlepage
\vspace{-0.35cm}
\begin{center}
\color{coursered}\bfseries Instructor material - do not distribute with the question deck
\end{center}
\end{frame}

\begin{frame}{Solution overview}
\begin{columns}[T,onlytextwidth]
\begin{column}{0.62\textwidth}
\Large This instructor deck provides a complete matched solution and grading guidance.

\vspace{0.35cm}
\normalsize
Coverage: Week 4: MDPs, Bellman backups, value iteration, and policy iteration
\end{column}
\begin{column}{0.33\textwidth}
\begin{block}{Points}10 total\end{block}
\begin{block}{Expected time}6-8 hours\end{block}
\begin{block}{Submission}Repository plus experiment memo\end{block}
\end{column}
\end{columns}
\end{frame}

\begin{frame}{Instructor verification checklist}
\small
\begin{itemize}
  \item Use Python 3.11 and NumPy only unless the prompt explicitly permits another package.
  \item Preserve the published interfaces and make all tie-breaking deterministic.
  \item State assumptions, validate inputs, and handle empty, terminal, and unreachable cases.
  \item Include focused tests whose names explain the defect each test would expose.
  \item Report measured evidence; do not replace experimental results with unsupported claims.
\end{itemize}
\end{frame}

\begin{frame}{Solution 1.a - MDP model (1 points)}
\small
\textbf{Question focus.} Define S, A, P, R, gamma, initial state, walls, and terminal states.

\vspace{0.25cm}
\answerlabel States are legal grid cells; actions are cardinal moves; intended motion has probability 0.8 and perpendicular slips 0.1 each. Blocked motion stays in place. Rewards and terminal semantics are stated, with gamma in [0,1).
\end{frame}

\begin{frame}{Solution 1.b - MDP model (1 points)}
\small
\textbf{Question focus.} Write validation tests for every state-action transition distribution.

\vspace{0.25cm}
\answerlabel For each legal nonterminal pair, successor probabilities are nonnegative and sum to one within tolerance; all successors are legal states. Terminal states expose no actions or a documented absorbing convention.
\end{frame}

\begin{frame}{Solution 2.a - Value iteration (1 points)}
\small
\textbf{Question focus.} Write the Bellman optimality backup used for each nonterminal state.

\vspace{0.25cm}
\answerlabel V\_new(s) is the maximum over actions of the sum over successors of P(s-prime|s,a) times [R(s,a,s-prime)+gamma V\_old(s-prime)]. Terminal values follow the declared reward convention.
\end{frame}

\begin{frame}{Solution 2.b - Value iteration (1 points)}
\small
\textbf{Question focus.} Use a copy of the previous value vector and stop when the maximum residual is below tolerance.

\vspace{0.25cm}
\answerlabel Synchronous updates read only V\_old. The residual is max\_s absolute(V\_new(s)-V\_old(s)); iteration count has a safety bound and tolerance is positive.
\end{frame}

\begin{frame}{Solution 2.c - Value iteration (1 points)}
\small
\textbf{Question focus.} Extract a deterministic greedy policy and report values for designated cells.

\vspace{0.25cm}
\answerlabel Each action is scored with the final Bellman lookahead; stable action order resolves ties. Reported cells match an independently computed oracle within tolerance.
\end{frame}

\begin{frame}{Solution 3.a - Policy iteration (1 points)}
\small
\textbf{Question focus.} Describe iterative policy evaluation for a fixed policy.

\vspace{0.25cm}
\answerlabel Repeatedly apply the Bellman expectation backup using only the current policy action until the evaluation residual is below tolerance.
\end{frame}

\begin{frame}{Solution 3.b - Policy iteration (1 points)}
\small
\textbf{Question focus.} Implement greedy improvement and state the stability condition.

\vspace{0.25cm}
\answerlabel For every nonterminal state choose an action maximizing one-step lookahead under the evaluated values. Stop when no chosen action changes under deterministic ties.
\end{frame}

\begin{frame}{Solution 3.c - Policy iteration (1 points)}
\small
\textbf{Question focus.} Compare the final policy and values with value iteration.

\vspace{0.25cm}
\answerlabel Optimal values agree within a declared numerical tolerance, and policies are equal or choose actions with equal optimal action value.
\end{frame}

\begin{frame}{Solution 4.a - Validation and analysis (1 points)}
\small
\textbf{Question focus.} Test gamma=0, a one-state terminal MDP, blocked slips, and a reward perturbation.

\vspace{0.25cm}
\answerlabel At gamma=0 only immediate expected reward matters. The terminal case is stable. Blocked slips accumulate probability on staying put. Increasing one terminal reward should not reduce optimal value in reachable states.
\end{frame}

\begin{frame}{Solution 4.b - Validation and analysis (1 points)}
\small
\textbf{Question focus.} Explain how gamma and slip probability change the policy in two experiments.

\vspace{0.25cm}
\answerlabel The answer compares controlled runs, reports values and actions, and links larger gamma to future rewards while larger slip can favor safer routes.
\end{frame}

\begin{frame}{Edge-case reasoning and expected behavior}
\small
\begin{itemize}
  \item Return the terminal convention immediately.
  \item Merge that probability into the stay-put successor.
  \item Use documented stable action order.
  \item Backups reduce to immediate expected reward.
  \item Reject the configuration before iteration.
\end{itemize}
\end{frame}

\begin{frame}{Grading rubric}
\scriptsize
\begin{tabularx}{\textwidth}{>{\bfseries}p{0.28\textwidth} p{0.08\textwidth} X}
\toprule
Category & Points & Full-credit evidence \\
\midrule
MDP model & 2 & Probabilities normalize and terminal/wall semantics are explicit. \\
Value iteration & 3 & Backups, convergence, and terminal handling are correct. \\
Policy iteration & 3 & Evaluation/improvement loops and stopping rule are sound. \\
Validation and analysis & 2 & Tests expose modeling errors and conclusions follow evidence. \\
\bottomrule
\end{tabularx}
\end{frame}

\begin{frame}{Reflection exemplar}
\small
\answerlabel A strong answer isolates a transition, reward, terminal, gamma, or slip convention, reports before-and-after values or policies, and explains why both solvers changed consistently.

\vspace{0.25cm}
\gradinglabel Award credit for a specific claim supported by technical evidence from the submitted work.
\end{frame}

\end{document}
